\documentclass[10pt,a4paper]{amsart}
\usepackage{amsmath,amssymb,amsthm,graphicx}
% Minimal class-compatibility preamble for a reader excerpt of
%   Frederic Deglise, Jean Fasel, Fangzhou Jin, Adeel A. Khan,
%   "On the rational motivic homotopy category",
%   Journal de l'Ecole polytechnique -- Mathematiques 8 (2021), DOI 10.5802/jep.153.
%
% The block below is the author's own macro, package and theorem declaration block of the
% pinned source file, copied mechanically (source lines 6-16 and 18-21 and 37-135). Only the
% publisher's OT2/cyrillic font block (source lines 22-35) is left out, because it declares
% 8-bit cyrillic fonts that the retained text never uses; a Unicode text font is declared
% instead so that the accented letters of the retained author text are rendered. No
% mathematical macro, statement, symbol or assumption is introduced, changed or removed.
\usepackage{fontspec}
\setmainfont{lmroman10-regular.otf}[
  ItalicFont=lmroman10-italic.otf,
  BoldFont=lmroman10-bold.otf,
  BoldItalicFont=lmroman10-bolditalic.otf]
\usepackage{xspace}

\def\To#1{\mathchoice{\xrightarrow{\textstyle\kern4pt#1\kern3pt}}{\stackrel{#1}{\longrightarrow}}{}{}}
\def\mto{\mathchoice{\longmapsto}{\mapsto}{\mapsto}{\mapsto}}
\newcommand{\Psfrac}[2]{(\sfrac{#1}{#2})}
\newcommand{\psfrac}[2]{\sfrac{(#1)}{#2}}
\newcommand{\spfrac}[2]{\sfrac{#1}{(#2)}}
\newcommand{\pspfrac}[2]{\sfrac{(#1)}{(#2)}}
\let\epsilon\varepsilon

\newcommand{\tstructure}{$t$-structure\xspace}
\newcommand{\tstructures}{$t$-structures\xspace}
\newcommand{\BM}{\mathrm{BM}}
\usepackage[all,2cell,emtex,cmtip]{xy}\let\labelstyle\textstyle
\usepackage{bbm,nicefrac}
\usepackage[mathscr]{euscript}
\usepackage{mathrsfs}
\newcommand{\modl}[1]{#1\textnormal{-}\mathrm{mod}}
\newcommand{\base}{\mathscr S}

\newcommand{\T}{\mathscr T}
\newcommand{\C}{\mathscr C}

\DeclareMathOperator{\uK}{\underline K}
\DeclareMathOperator{\uPic}{\underline{Pic}}
\DeclareMathOperator{\uZZ}{\underline{\ZZ}}
\DeclareMathOperator{\rk}{rk}
\DeclareMathOperator{\GW}{GW}

\renewcommand{\AA}{\mathbb A}
\newcommand{\E}{\mathbb E}
\newcommand{\F}{\mathbb F}
\newcommand{\un}{\mathbbm 1}
\newcommand{\MSL}{\mathbf{MSL}}
\newcommand{\KQ}{\mathbf{KQ}}
\newcommand{\KW}{\mathbf{KW}}
\newcommand{\KMW}{\mathbf K^\mathrm{MW}}
\newcommand{\uW}{\underline{\mathrm W}}
\newcommand{\HuW}{\mathbf H \underline{\mathrm W}}
\newcommand{\HM}{\mathbf H_{\mathrm{M}}}
\newcommand{\HMW}{\mathbf H_{\mathrm{MW}}}
\newcommand{\HW}{\mathbf H_{\mathrm{W}}}
\newcommand{\HH}{\mathrm H}
\newcommand{\HA}{\mathrm H_{\AA^1}}
\newcommand{\BA}{\mathrm H^{\AA^1}}

\newcommand{\GL}{\mathrm{GL}}
\newcommand{\SL}{\mathrm{SL}}

\newcommand{\thom}{\mathfrak t}
\newcommand{\vb}[1]{\langle#1\rangle}

\DeclareMathOperator{\codim}{codim}
\newcommand{\eff}{\mathrm{eff}}

\newcommand{\sm}{\mathrm{Sm}}
\DeclareMathOperator{\DM}{DM}
\DeclareMathOperator{\DMB}{DM_\mathcyr B}
\DeclareMathOperator{\SH}{SH}
\DeclareMathOperator{\DA}{D_{\AA^1}}
\DeclareMathOperator{\DAeff}{D_{\AA^1}^{\eff}}
\DeclareMathOperator{\DAminus}{D_{\AA^1,-}}
\DeclareMathOperator{\Id}{Id}
\DeclareMathOperator{\Hom}{Hom}
\DeclareMathOperator{\Maps}{Maps}
\DeclareMathOperator{\uHom}{\underline{Hom}}
\DeclareMathOperator{\Pic}{Pic}
\DeclareMathOperator{\CH}{CH}
\DeclareMathOperator{\wCH}{\widetilde{CH}}
\DeclareMathOperator{\spec}{Spec}
\DeclareMathOperator{\Ind}{Ind}
\DeclareMathOperator{\ilim}{\varinjlim}
\DeclareMathOperator{\plim}{\varprojlim}
\newcommand{\twist}[1]{\langle #1 \rangle}
\DeclareMathOperator{\Th}{Th}
\DeclareMathOperator{\Tw}{Tw}
\DeclareMathOperator{\tdet}{\widetilde \det}
\DeclareMathOperator{\tomega}{\widetilde \omega}
\DeclareMathOperator{\K}{K}
\DeclareMathOperator{\Der}{D}
\DeclareMathOperator{\Spt}{Spt}

\newcommand{\ZZ} {\mathbb Z}
\newcommand{\ZZot} {\mathbb Z[\nicefrac 1 2]}
\newcommand{\ot} {{\nicefrac 1 2}}
\newcommand{\QQ} {\mathbb Q}
\newcommand{\OO}{\mathcal O}
\newcommand{\PP} {\mathbb P}
\newcommand{\GG} {\mathbb{G}_m}
\newcommand{\GGx}[1] {\mathbb{G}_{m,#1}}
\newcommand{\nis}{\mathrm{Nis}}
\newcommand{\zar}{\mathrm{Zar}}

\newtheorem{thm}{Theorem}[section]
\newtheorem{prop}[thm]{Proposition}
\newtheorem{lm}[thm]{Lemma}
\newtheorem{cor}[thm]{Corollary}
\newtheorem{thmi}{Theorem}
\newtheorem{cori}[thmi]{Corollary}
\newtheorem*{lmi}{Key Lemma}

\renewcommand{\thethmi}{\Alph{thmi}}

\theoremstyle{remark}
\newtheorem{rem}[thm]{Remark}
\newtheorem{notation}[thm]{Notation}
\newtheorem{ex}[thm]{Example}
\newtheorem{conv}[thm]{Convention}
\newtheorem{setup}[thm]{Setup}
\newtheorem{warn}[thm]{Warning}

\theoremstyle{definition}
\newtheorem{df}[thm]{Definition}

\let\thesubsubsection\thethm
\newcommand{\numer}[1]{\refstepcounter{thm}\subsubsection{#1}}
\newtheorem{num}[thm]{}
\numberwithin{equation}{thm}

\renewcommand{\theequation}{\thethm.\alph{equation}}

% Publisher-class abbreviations (cedram) used by the retained author text.
\providecommand{\loccit}{\emph{loc. cit.}}
\providecommand{\bibliofont}{\small}
\providecommand{\backmatter}{}
\providecommand{\skpt}{\smallskip}
\providecommand{\smaller}{\small}
\newcommand{\ie}{i.e.,\xspace}
\newcommand{\eg}{e.g.,\xspace}
\newcommand{\etc}{etc.\xspace}
\newcommand{\cf}{cf.\xspace}
\newcommand{\resp}{resp.\xspace}
\usepackage{cleveref}
\crefname{originallemma}{original lemma}{original lemmas}
\Crefname{originallemma}{Original lemma}{Original lemmas}
\makeatletter
\newlabel{lm:fiber_Q_eq}{{2.2}{source}{}{sourcepointer-1}{}}
\newlabel{lm:fiber_Q_eq@cref}{{[originallemma][0][]2.2}{[1][0][]source}{}{}{}}
\makeatother
\title{The 2-inverted Witt sheaf of a Z[1/2]-scheme is the pushforward of the Witt sheaf of its characteristic-zero fibre}
\author{Frédéric Déglise, Jean Fasel, Fangzhou Jin, Adeel A. Khan}
\date{doi:10.5802/jep.153; CC BY 4.0}
\begin{document}
\maketitle
\setcounter{section}{13}
\setcounter{thm}{1}
\noindent\textit{Editorial note.} The counted unit is Proposition 13.2 at source lines 2513-2519, source label prop:Witt\_2inverted, with the proof environment that begins on the next line at 2520-2568; the pairing is the ordinary adjacent one. The counted statement and its proof are placed immediately after the editorial material and the printed number is fixed by initial\_counters so that the proposition carries its original number 13.2; the definitional passages retained after it print with excerpt-local numbers because the items that separate them in the published article are not retained (their original numbers are 2.1, 5.1, 5.2 and 13.1).\par\medskip
\noindent\textit{Editorial note.} The appendix opening at source lines 2496-2512 is retained because the proof uses Balmer's higher Witt groups with support introduced there; the definition of the Witt sheaf at 891-894, the characteristic-zero fibre paragraph at 614-621 and the comparison map at 906-916 are retained because the statement names them. These four context passages follow the counted unit in the excerpt; none of them was altered.\par\medskip
\noindent\textit{Editorial note.} One label is outside the retained spans: lm:fiber\_Q\_eq, which occurs only in the sentence at source line 2511 that compares the counted proposition with the paper's key lemma. It is declared as a non-load-bearing external reference with its original number 2.2 and its source position, so the material carries a pointer rather than a replacement proof.\par\medskip
\noindent\textit{Editorial note.} The retained preamble reproduces the author's macro, package and theorem declaration block verbatim from source lines 6-16, 18-21 and 37-135. The publisher font block at source lines 22-35 is left out because it declares 8-bit cyrillic fonts that the retained text never uses; a Unicode text font is declared instead so that the accented letters inside the retained author sentences are rendered. The cedram publisher class is not installed on the assembly host, so the wrapper class amsart is used; no mathematical macro, symbol, statement or hypothesis is added, changed or removed.\par\medskip
\noindent\textit{Editorial note.} The retained bibliography carries exactly the five works cited by the retained text (Balmer, Witt cohomology; Balmer, Witt groups; Bondarko-Deglise; Balmer-Walter; Lam), copied verbatim from the public corpus record of this same source version.\par\medskip
\noindent\textit{Editorial scope.} One original result: Proposition 13.2 (label prop:Witt\_2inverted) of the source with its complete author proof at source lines 2520-2568, together with the appendix definitions and the comparison map that the statement names. The already-counted Theorem 3.6 clause 2 of the same source, the Witt-motives identification and the Gersten-Witt corollary are not counted here and are not used as premises.\par\medskip
\noindent\textit{Allowed standard prerequisites (editorial).}\begin{itemize}
\item Witt groups and Witt sheaves of schemes. For a scheme S the Witt group W(S) is the monoid of isometry classes of symmetric bundles modulo metabolic ones, and the Witt sheaf on the site of smooth S-schemes is the Zariski sheaf associated with the presheaf X maps to W(X); the Lambda-linear version tensors with a subring Lambda of Q. The authors recall this at source lines 878-894 and both objects are retained in the material.. Fixes both sides of the counted isomorphism: the Lambda-linear Witt sheaf of S and the pushforward of the Lambda-linear Witt sheaf of the characteristic-zero fibre S\_Q.
\item Balmer's Witt groups with support and their localization sequence. For a closed pair (X,Z) there are triangulated Witt groups W\textasciicircum{}i(X,Z) fitting in the long exact localization sequence of Balmer-Witt theory, and for a closed subscheme Z of Spec Z with open complement U the sequence W\textasciicircum{}0(S,S\_Z) to W\textasciicircum{}0(S) to W\textasciicircum{}0(S\_U) to W\textasciicircum{}1(S,S\_Z) is exact; the authors recall this at source lines 2499-2512 and cite Balmer, Witt cohomology (1.7 and Th. 1.6).. Produces the exact sequence whose middle map is the map of the counted proposition and reduces the claim to the vanishing of W\textasciicircum{}i(S,S\_p)[1/2] for odd primes p at source lines 2529-2544.
\item Gersten-Witt spectral sequence and purity for Witt groups. For a regular scheme S with the dimension function delta = -codim there is a converging spectral sequence with E\textasciicircum{}1 terms W\textasciicircum{}\{-i-j\}(S\_(x), \{x\}) over the points of prescribed delta-value, and purity (devissage) identifies W\textasciicircum{}\{-i-j\}(S\_(x),\{x\}) with W(kappa(x)) when i+j is 0 mod 4 and with 0 otherwise; the authors cite Balmer-Walter (Th. 6.1) at source lines 2546-2564.. Computes the E\textasciicircum{}1 page whose terms are 2-primary torsion in positive characteristic, which is the observation that finishes the counted proof.
\end{itemize}\medskip
\ifdefined\hypertarget\hypertarget{sourcepointer-1}{}\fi
\noindent\textit{Source reference.} Original lemma 2.2: sources/algebra-fl04-on-the-rational-motivic-homotopy-c/pinned-original-source.tex, source lines 623--632. Only the sentence at source line 2511 compares the counted proposition with the paper's key lemma; the retained proof at 2520-2568 is a direct argument about Witt groups with support and never uses that lemma, so the pointer is navigation only.\par\medskip
\begin{prop}\label{prop:Witt_2inverted}
Let $S$ be a $\ZZot$-scheme and $\Lambda \subseteq \QQ$ a subring with $2 \in \Lambda^\times$.
Let $\nu:S_\QQ \to S$ be the pro-open immersion of the characteristic $0$ part of $S$
as in Paragraph \ref{num:zero_part}.

Then the map $\uW_S^\Lambda \to \nu_*(\uW_{S_\QQ}^\Lambda)$ from \eqref{eq:funct1_uW} is an isomorphism.
\end{prop}

\begin{proof}
It will be sufficient to prove that for any regular scheme $S$, the following map on Witt groups
is an isomorphism:
\[
W(S)[2^{-1}] \To{\nu^*} W(S_\QQ)[2^{-1}].
\]
We will use Balmer-Witt groups with support, as defined in \cite[\S 1]{BalmerWitt}.
Given a closed pair $(X,Z)$, we write $W^i(X,Z)$ for the $i$-th Witt group of $X$ with support in $Z$.

Consider a closed subscheme $Z$ of $\spec(\ZZ)$, with open complement $U$.
Write $S_Z$, $S_U$ for the pullbacks of $S/\spec(\ZZ)$ to $Z$, $U$.
Using the long exact sequence of Witt groups with support \cite[1.7]{BalmerWitt},
we get an exact sequence:
\[
W^0(S,S_Z) \to W^0(S) \to W^0(S_U) \to W^1(S,S_Z).
\]
Taking the filtered inductive limit over the category of closed subschemes of $S$,
and using additivity in $S_Z$, we get an exact sequence:
\[
\oplus_{p\neq 2} W^0(S,S_p) \to W^0(S) \To{\nu^*} W^0(S_\QQ) \to \oplus_{p\neq 2} W^1(S,S_p),
\]
where $p$ runs over the odd primes.

Therefore, to get the lemma, it is sufficient to prove that for any odd prime $p$,
the triangular Witt group $W^i(S,S_p)[2^{-1}]$ is zero.

Let $\delta$ be the dimension function on $S$ such that $\delta=-\codim_S$
(see \eg \cite[\S 1.1]{BD1}).
From the long exact sequence of Witt groups with support (op.\,cit., Th.\,1.6),
one can build out of the $\delta$-niveau filtration on $S_p$,
as in \cite[\S 3.1]{BD1}, and deduce a converging spectral
sequence\footnote{This construction is nothing else than the version with support of the Gersten--Witt spectral sequence,
whose first version can be found in \cite{BalmerWalter}.} of the form:
\[
E^1_{i,j}=\oplus_{x \in S_{p,(i)}} W^{-i-j}(S_{(x)},\{x\}) \implies W^{-i-j}(S,S_p),
\]
where $S_{p,{(i)}}$ denotes the set of points of $x$ such that $\delta(x)=i$.
Since $S$ is regular and defined over $\ZZot$, purity (dévissage) for the Witt groups \cite[Th.\,6.1]{BalmerWalter} yields:
\[
W^{-i-j}(S_{(x)},\{x\})=
\begin{cases}
W(\kappa(x)) & i+j\equiv 0 \bmod 4 \\
0 & \text{otherwise.}
\end{cases}
\]
This group is $2$-primary torsion as $\kappa(x)$ has positive characteristic
(\cite[VIII\,3.2 \& 6.15]{LamQuadratic}).
In particular $E^1_{i,j}[\ot]$ vanishes and one can conclude.
\end{proof}

\begin{num}\label{num:zero_part}
For a scheme $X$, denote by $X_\QQ = X\times_{\spec(\ZZ)}\spec(\QQ)$ the characteristic zero fiber.
We denote the inclusion by
\begin{equation*}
\nu = \nu_X : X_\QQ \to X.
\end{equation*}
The following observation is the \textit{key lemma} of this paper:
\end{num}

\begin{df}\label{df:unramified_Witt}
For any scheme $S$, we define the \emph{Witt sheaf} $\uW_S$ on the site of smooth $S$-schemes as the Zariski sheaf associated with the presheaf $X \mto W(X)$.
Given a subring $\Lambda$ of $\QQ$, we set $\uW_S^\Lambda=\uW_S \otimes_\ZZ \Lambda$.
\end{df}

\begin{rem}
Note that $\uW_S$ can also be described as the restriction to $\sm_S$ of the sheaf $\uW$ associated with $W$ on the big Zariski site of schemes of finite presentation over $S$.
In particular, for any morphism of schemes $f:X \to S$ one gets a morphism of sheaves:
\begin{equation}\label{eq:funct1_uW}
\uW_S \to f_*(\uW_X),
\end{equation}
or equivalently by adjunction:
\begin{equation}\label{eq:funct2_uW}
f^*(\uW_S) \to \uW_X.
\end{equation}
\end{rem}

This appendix brings some complements to the theory of Witt and unramified Witt sheaves
(see Definition~\ref{df:unramified_Witt}) of independent interest.

\begin{num}\label{num:Balmer}
The theory of Witt groups has been extended to
higher Witt groups $W^i(S)$ by Balmer (see \eg \cite{BalmerBasic}).
Recall that these groups are defined by considering the deri\-ved category
$\T=\Der^b(\mathrm{VB}_S)$ of the exact category of vector bundles over $S$
equipped with the $(-1)^i$-duality $D_i=\uHom(-,\mathcal O_S)[i]$
(see \cite[1.4.12]{BalmerBasic}).
It is then possible to mimic the classical definition of
symmetric bilinear forms/symmetric objects, with respect to the duality $D_i$, and metabolic ones (called ``neutral'' in loc.\,cit.).
One defines $W^i(X)$ as the monoid of isomorphisms classes of
$D_i$-symmetric objects of $\T$ modulo the submonoid of $D_i$-metabolic ones (see loc.\,cit. 1.4.3).

The following lemma is the Witt-version of our \emph{key lemma} (Lemma~\ref{lm:fiber_Q_eq}).
\end{num}

\begin{thebibliography}{99}
\bibitem{BalmerWitt} P. Balmer - “Witt cohomology, Mayer-Vietoris, homotopy invariance and the Gersten conjecture” , $K$-Theory 23 (2001) no. 1, p. 15-30
\bibitem{BalmerBasic} P. Balmer - “Witt groups” , in Handbook of $K$-theory. Vol. 1, 2 , Springer, Berlin, 2005, p. 539-576
\bibitem{BD1} M. Bondarko \& F. Déglise - “Dimensional homotopy t-structures in motivic homotopy theory” , Adv. Math. 311 (2017), p. 91-189
\bibitem{BalmerWalter} P. Balmer \& C. Walter - “A Gersten-Witt spectral sequence for regular schemes” , Ann. Sci. École Norm. Sup. (4) 35 (2002) no. 1, p. 127-152
\bibitem{LamQuadratic} T. Y. Lam - Introduction to quadratic forms over fields , Graduate Studies in Math. , vol. 67 , American Mathematical Society, Providence, RI, 2005
\end{thebibliography}
\end{document}
